\documentclass[11pt,a4paper]{article}

\usepackage[margin=2.4cm]{geometry}
\usepackage{graphicx}
\usepackage{booktabs}
\usepackage{amsmath}
\usepackage{siunitx}
\usepackage[colorlinks=true,linkcolor=black,urlcolor=blue,citecolor=black]{hyperref}
\usepackage{caption}
\captionsetup{font=small,labelfont=bf}

\sisetup{detect-all}

\title{Arm azimuth and corner sharpness in quadrotor gate racing}
\author{ARIEL / SPEAR}
\date{2026-08-17}

% Generated by Reports/make_figures.py -- do not edit.
\newcommand{\FeasLo}{20.44}
\newcommand{\FeasHi}{69.56}
\newcommand{\NPoints}{7}
\newcommand{\RollNarrow}{313.3}
\newcommand{\RollWide}{121.3}
\newcommand{\RollRatio}{2.58}
\newcommand{\PitchNarrow}{120.8}
\newcommand{\TorqueLo}{1.489}
\newcommand{\TorqueHi}{3.996}
\newcommand{\XTrackSixty}{0.55}
\newcommand{\XTrackNinety}{0.61}
\newcommand{\XTrackOneTwenty}{0.47}
\newcommand{\XTrackMeanLo}{0.24}
\newcommand{\XTrackMeanHi}{0.35}
\newcommand{\PubBestSixty}{12.172}
\newcommand{\PubBestNinety}{8.578}
\newcommand{\PubBestOneTwenty}{7.078}
\newcommand{\PrevBestSixty}{10.164}
\newcommand{\PrevBestNinety}{8.367}
\newcommand{\PrevBestOneTwenty}{6.648}
\newcommand{\BestSixty}{10.398}
\newcommand{\ArgmaxSixty}{36.81}
\newcommand{\NTiesSixty}{1}
\newcommand{\SpreadSixty}{4.1}
\newcommand{\NarrowSixty}{-3.0}
\newcommand{\WideSixty}{-3.8}
\newcommand{\BestNinety}{8.250}
\newcommand{\ArgmaxNinety}{45.00\text{--}69.56}
\newcommand{\NTiesNinety}{2}
\newcommand{\SpreadNinety}{32.7}
\newcommand{\NarrowNinety}{-32.7}
\newcommand{\WideNinety}{-0.5}
\newcommand{\BestOneTwenty}{7.117}
\newcommand{\ArgmaxOneTwenty}{53.19\text{--}69.56}
\newcommand{\NTiesOneTwenty}{3}
\newcommand{\SpreadOneTwenty}{25.8}
\newcommand{\NarrowOneTwenty}{-25.8}
\newcommand{\WideOneTwenty}{0.0}
\newcommand{\GateSpeedSixty}{5.60}
\newcommand{\GateLoSixty}{3.84}
\newcommand{\GateHiSixty}{6.60}
\newcommand{\MeanSpeedSixty}{5.89}
\newcommand{\StartupFacSixty}{1.77}
\newcommand{\GearingFacSixty}{1.05}
\newcommand{\SpeedRatioSixty}{1.86}
\newcommand{\Polyline}{38.4}
\newcommand{\PeakRef}{29}
\newcommand{\GateSpeedNinety}{3.86}
\newcommand{\GateLoNinety}{3.08}
\newcommand{\GateHiNinety}{5.34}
\newcommand{\MeanSpeedNinety}{5.22}
\newcommand{\StartupFacNinety}{1.58}
\newcommand{\GearingFacNinety}{1.35}
\newcommand{\SpeedRatioNinety}{2.14}
\newcommand{\GateSpeedOneTwenty}{2.44}
\newcommand{\GateLoOneTwenty}{2.11}
\newcommand{\GateHiOneTwenty}{4.55}
\newcommand{\MeanSpeedOneTwenty}{4.77}
\newcommand{\StartupFacOneTwenty}{1.49}
\newcommand{\GearingFacOneTwenty}{1.96}
\newcommand{\SpeedRatioOneTwenty}{2.92}
\newcommand{\ArcLo}{39.4}
\newcommand{\ArcHi}{40.0}
\newcommand{\GateSpreadSixty}{3.0}
\newcommand{\GateSpreadNinety}{26.8}
\newcommand{\GateSpreadOneTwenty}{21.4}
\newcommand{\GearEnd}{4.99}
\newcommand{\GearInterior}{2.50}
\newcommand{\GearRatio}{2.00}


\begin{document}
\maketitle

\begin{abstract}
%A quadrotor's four arms are swept through the full range of azimuths that keeps its propellers clear of one another, and each resulting airframe is scored by the fastest speed at which it still completes a slalom, where completing means passing inside every gate opening. Three corner sharpnesses are flown. No single geometry is best across them: the optimum migrates monotonically toward narrower higher-roll-authority frames as corners sharpen, and the penalty for being wrong grows with them---from \SI{\SpreadSixty}{\percent} across the family at $60^\circ$ to \SI{\SpreadOneTwenty}{\percent} at $120^\circ$, where the narrowest frame is outright best and the widest last. The effect is visible only once the controller cascade is tuned so that the airframe, rather than the control loop, is the binding constraint.
This study supports a potential benefit of actively morphed drones. We show that certain morphologies are better suited to a given task than others. We fly a quadrotor with varying arm azimuth angles through three slalom courses, each with a different corner sharpness. The performance of a given morphology is evaluated by the maximum speed at which it can complete the course without missing a gate. The best morphology (i.e., the arm azimuth angle) depended on the corner sharpness. This justifies the need for actively morphed drones in gate-passing tasks.
\end{abstract}


\section*{Notation}

Two symbols are reused in the drone-racing literature and are kept distinct
here: $t$ is always the arm half-angle, never time, and $L$ is always arm
length, never path length. Since $t$ is unavailable, time is written $T$ where
it is named, and an overdot denotes a rate of change with respect to time.

\begin{center}
\small
\begin{tabular}{@{}l l l@{}}
  \toprule
  \multicolumn{3}{@{}l}{\textit{Airframe}} \\
  $t$ & arm half-angle, the swept parameter & \si{\degree} \\
  $L$ & arm length, hub to rotor centre & \si{\metre} \\
  $r$ & propeller radius & \si{\metre} \\
  $\tau$ & propeller clearance tolerance, \num{0.1} & --- \\
  $I$ & body inertia tensor about the centre of gravity & \si{\kilo\gram\metre\squared} \\
  $I_{ij}$ & its $(i,j)$ entry; $i \neq j$ are the products of inertia & \si{\kilo\gram\metre\squared} \\
  $B_M$ & moment allocation matrix, column $i$ the moment per newton of rotor $i$ & \si{\metre} \\
  $\alpha_\phi,\ \alpha_\theta,\ \alpha_\psi$ & roll, pitch and yaw
    angular-acceleration authority, the row norms of $I^{-1}B_M$ &
    \si{\radian\per\second\squared\per\newton} \\
  $\psi$ & heading (yaw), commanded to follow the path tangent & \si{\degree} \\
  \midrule
  \multicolumn{3}{@{}l}{\textit{Course}} \\
  $\theta$ & turn angle demanded at each gate: the difficulty knob & \si{\degree} \\
  $s$ & leg length, gate to gate along the path; held fixed at \SI{2.4}{\metre} & \si{\metre} \\
  $d$ & gate spacing along the course axis & \si{\metre} \\
  $A$ & lateral amplitude of the weave & \si{\metre} \\
  $R$ & corner radius & \si{\metre} \\
  $a_{\mathrm{lat}}$ & lateral acceleration demanded at speed $v$ & \si{\metre\per\second\squared} \\
  $\ell$ & total path length, lead-out included & \si{\metre} \\
  $v$ & nominal speed; see the note on what it measures & \si{\metre\per\second} \\
  \midrule
  \multicolumn{3}{@{}l}{\textit{Trajectory and controller}} \\
  $P(u)$ & the reference B-spline: position as a function of its parameter & \si{\metre} \\
  $u$ & spline parameter; the curve advances at constant $\dot u$ while cruising & --- \\
  $T_{\text{startup}}$ & duration of the ramp from rest, \SI{3}{\second} & \si{\second} \\
  $T_{\text{total}}$ & time allotted for the whole course & \si{\second} \\
  $\omega_n$ & commanded closed-loop natural frequency (bandwidth scale) & \si{\radian\per\second} \\
   & its $-3$\,dB bandwidth is $0.644\,\omega_n$ at $\zeta=1$ & \\
  $\zeta$ & commanded closed-loop damping ratio; \num{1} is critical damping & --- \\
  $\Omega$ & body angular rate & \si{\radian\per\second} \\
  $K_R,\ K_\Omega$ & attitude proportional and rate gains, $K_R=\omega_n^2 I$ and
    $K_\Omega = 2\zeta\omega_n I$ & --- \\
  $m,\ g$ & total mass and gravitational acceleration & \si{\kilogram},
    \si{\metre\per\second\squared} \\
  $p$ & measured position & \si{\metre} \\
  $\dot p$ & measured velocity & \si{\metre\per\second} \\
  $p_{\rm ref}$ & the reference's position & \si{\metre} \\
  $\dot p_{\rm ref}$ & the reference's velocity & \si{\metre\per\second} \\
  $\ddot p_{\rm ref}$ & the reference's acceleration & \si{\metre\per\second\squared} \\
  $a_{\rm cmd}$ & acceleration commanded by the position loop
    (Appendix~\ref{sec:feedforward}) & \si{\metre\per\second\squared} \\
  $F$ & commanded force, $m(a_{\rm cmd} - g)$ & \si{\newton} \\
  $K_p,\ K_v$ & position and velocity feedback gains & --- \\
  $x,\ y,\ z$ & body-fixed axes, $x$ forward, $y$ right, $z$ down & \si{\metre} \\
  \bottomrule
\end{tabular}
\end{center}

\section{Airframes}

% Every body is the same quadrotor---\SI{0.829}{\kilo\gram}, \SI{5}{inch}
% propellers ($r=\SI{0.0635}{\metre}$), \SI{0.20}{\metre} arms, thrust-to-weight
% \num{5.24}, matching the SPEAR reference airframe---differing only in the
% azimuth of its arms. With the arms at $\pm t$ and $180^\circ \pm t$ the family is
% bilaterally symmetric and spans one parameter. Propeller--propeller clearance
% bounds it: neighbouring rotors are $2L\sin t$ and $2L\cos t$ apart and both must
% exceed $2(1+\tau)r$ with a tolerance $\tau=0.1$, giving
% $t \in [\FeasLo^\circ, \FeasHi^\circ]$, sampled at \NPoints{} points
% (Figure~\ref{fig:morph}).
Consider a quadrotor with \SI{0.829}{\kilo\gram} mass, \SI{5}{inch} propellers ($r=\SI{0.0635}{\metre}$), \SI{0.20}{\metre} arms, and thrust-to-weight ratio of \num{5.24}. We vary the azimuth of its arms, parameterized by the half-angle $t$,  measured from the body-fixed longitudinal axis $x$. We maintain the bilaterally symmetric X configuration by placing the arms at $\pm t$ and $180^\circ \pm t$. Inter-propeller clearance is maintained by considering configurations in $t \in [\FeasLo^\circ, \FeasHi^\circ]$, sampled at \NPoints{} points (Figure~\ref{fig:morph}).
\begin{figure}[htb]
  \centering
  \includegraphics[width=0.78\textwidth]{figures/morphologies.pdf}
  \caption{The seven airframes, drawn to scale; the arrow marks the x-axis and the direction of travel and the blue disc the \SI{0.05}{\metre} body. The drone is commanded to point along its path so the weave is flown by banking. Thus, the roll authority is the discriminating quantity. At low $t$, the frame stretches fore-and-aft. The classic X quad is at $t = 45^\circ$. At high $t$, the frame is laterally stretched. Beneath each, the roll and pitch angular-acceleration authorities $\alpha_\phi$, $\alpha_\theta$ in \si{\radian\per\second\squared} per newton of thrust, defined as the row norms of $I^{-1}B_M$. As the $t$ increases, the roll and pitch authorities swap their relative importance. The slalom, however demands roll authority.}
  \label{fig:morph}
\end{figure}

% Roll authority $\alpha_\phi$ varies by a factor of \num{\RollRatio} across the family, from \num{\RollNarrow} at the narrow end to \num{\RollWide} at the wide end. It is deliberately inertia-normalised: raw torque authority gives the \emph{opposite} ordering, since widening the frame lengthens the lever arms (\SIrange{\TorqueLo}{\TorqueHi}{\newton\metre}) but raises inertia faster.
Roll authority $\alpha_\phi$ varies by a factor of \num{\RollRatio} across the studied configurations, from \num{\RollNarrow} at the narrow end to \num{\RollWide} at the wide end. The Roll Authority is deliberately inertia-normalised: raw torque authority gives the \emph{opposite} ordering, since widening the airframe lengthens the lever arms (\SIrange{\TorqueLo}{\TorqueHi}{\newton\metre}) but raises inertia faster.

\section{Courses}

% Difficulty is one number, the turn angle $\theta$ demanded at each gate, at fixed
% leg length $s=\SI{2.4}{\metre}$ (Figure~\ref{fig:courses}). Consecutive legs run
% at $\pm\theta/2$ to the course axis, so
% \begin{equation}
%   d = s\cos(\theta/2), \qquad
%   A = \tfrac{s}{2}\sin(\theta/2), \qquad
%   R = \frac{s}{2\sin(\theta/2)}, \qquad
%   a_{\mathrm{lat}} = \frac{v^2}{R},
% \end{equation}
% for gate spacing, lateral amplitude, corner radius and the lateral acceleration
% demanded at speed $v$. Holding $s$ fixed rather than $d$ is what makes $R$ fall
% monotonically with $\theta$: at fixed spacing a sharper turn also lengthens the
% legs, and the corner ends up \emph{wider}.

% \begin{figure}[htb]
%   \centering
%   \includegraphics[width=0.92\textwidth]{figures/courses.pdf}
%   \caption{The three courses. Red bars are the fifteen scored gates
%   (\SI{1.0}{\metre} opening), grey the two lead-out waypoints that let the drone
%   fly \emph{through} the finish rather than stop on it, and the green dot the
%   spawn. Gate planes face along the local direction of travel, which is what the
%   pass detector requires.}
%   \label{fig:courses}
% \end{figure}
Difficulty is given by the turn angle $\theta$ demanded at each gate, at fixed leg length $s=\SI{2.4}{\metre}$ (Figure~\ref{fig:courses}). Consecutive legs run at $\pm\theta/2$ to the course axis, so

\begin{equation}
  d = s\cos(\theta/2), \qquad
  A = \tfrac{s}{2}\sin(\theta/2), \qquad
  R = \frac{s}{2\sin(\theta/2)}, \qquad
  a_{\mathrm{lat}} = \frac{v^2}{R},
\end{equation}

for gate spacing, lateral amplitude, corner radius, and the lateral acceleration demanded at speed $v$. Holding $s$ fixed rather than $d$ is what makes $R$ fall monotonically with $\theta$: at fixed spacing, a sharper turn also lengthens the legs, and the corner ends up \emph{wider}.



\begin{figure}[htb]

  \centering

  \includegraphics[width=0.92\textwidth]{figures/courses.pdf}

  \caption{The three courses. Red bars are the fifteen scored gates  (\SI{0.5}{\metre} radius), grey bars are the two lead-out waypoints that let the drone fly \emph{through} the finish rather than stop on it, and the green dots are the spawn points. Gate planes face along the local direction of travel, which is what the pass detector requires.}

  \label{fig:courses}

\end{figure}

\section{Scoring}

% Each airframe is scored by the fastest speed at which it still completes the
% course, located by bisection. Fixed-speed scoring was abandoned: across nine
% operating points from \SIrange{2.0}{3.8}{\metre\per\second} every airframe scored
% identically, and a \SI{5}{\percent} increase collapsed them all. Maximum
% completing speed is continuous by construction and needs no operating point to be
% calibrated. ``Finished'' has to be defined carefully, though: a gate counts as
% passed only if the drone's offset from the gate centre, measured \emph{in the
% gate plane at the crossing} and including its vertical component, is inside the
% half-opening. Two weaker tests were in use earlier and both inflate the result;
% see the note below. Bisection is valid because completion is monotone in speed; the
% assumption is checked on every run and held in all cases reported here.

% Two conditions matter for the numbers below. Feedforward of the reference
% velocity and acceleration is enabled, and the controller cascade is tuned so the
% position loop sits well below the attitude loop
% ($\omega_n^{\mathrm{att}}=\SI{24}{\radian\per\second}$,
% $\omega_n^{\mathrm{pos}}=\SI{2.0}{\radian\per\second}$, $\zeta=1$). Under the
% inherited tuning the loops were separated by only a factor of \num{3.2}, motor
% saturation sat at \SIrange{0.1}{2}{\percent}, and the \emph{controller} was the
% binding constraint---in that regime the same airframe wins on every course and
% the spread across the family is \SI{3}{\percent}. At the tuned operating point
% every airframe saturates its motors \SIrange{15}{21}{\percent} of the time, so
% the airframe binds instead.

% \paragraph{Why one tuning is fair across airframes.} The gains are not shared
% between bodies; the \emph{closed-loop response} is. Attitude gains are rederived
% per morphology as $K_R=\omega_n^2 I$ and $K_\Omega=2\zeta\omega_n I$, so $I$
% cancels and every airframe is commanded to the same
% $(\omega_n,\zeta)=(\SI{24}{\radian\per\second},1)$ regardless of its inertia.
% The gain \emph{magnitudes} therefore differ between bodies, and between axes of
% one body, in proportion to the corresponding moments. This cancellation is exact
% only when the body axes are principal axes; for this family the products of
% inertia are identically zero (measured $\max|I_{ij}|=0$ for $i\neq j$, not merely
% small), because it retains both mirror planes. What is compared is thus geometry
% at matched commanded bandwidth, not geometry plus an accidental tuning
% advantage.
Each airframe is scored by the maximum speed at which it still completes the course without missing a gate. The speeds are located by bisection. A gate counts as passed only if the drone's offset from the gate centre, measured \emph{in the gate plane at the crossing} and including its vertical component, is within the gate radius of \SI{0.5}{\metre}. That offset is a Euclidean norm, so the gate is a circular opening rather than the square frame the word ``gate'' usually denotes; a square of the same half-width encloses \SI{27}{\percent} more area, all of it near the diagonals, so the criterion used here is the stricter of the two. Bisection is valid because completion is monotone in speed; the assumption is checked on every run and held in all cases reported here.

Two conditions matter for the numbers below. Feedforward of the reference velocity and acceleration is enabled (c.f. section \ref{sec:feedforward}), and the controller cascade (outer position loop and inner attitude loop) is tuned so the position loop sits well below the attitude loop ($\omega_n^{\mathrm{att}}=\SI{24}{\radian\per\second}$, $\omega_n^{\mathrm{pos}}=\SI{2.0}{\radian\per\second}$, $\zeta=1$). Under the inherited tuning, the loops were separated by only a factor of \num{3.2}, motor saturation sat at \SIrange{0.1}{2}{\percent}, and the \emph{controller} was the active constraint. In that regime, the same airframe wins on every course, and the spread across the family is \SI{3}{\percent}. At the tuned operating point, every airframe saturates its motors \SIrange{15}{21}{\percent} of the time, so the airframe binds the performance instead.

\paragraph{Why one tuning is fair across airframes.} The gains are not shared between bodies; the \emph{closed-loop response} for set point hover is. Attitude gains are rederived per morphology as $K_R=\omega_n^2 I$ and $K_\Omega=2\zeta\omega_n I$, so $I$ cancels and every airframe is commanded to the same $(\omega_n,\zeta)=(\SI{24}{\radian\per\second},1)$ regardless of its inertia. The gain \emph{magnitudes} therefore differ between bodies, and between body axes, in proportion to the corresponding moments. This cancellation is exact only when the body axes are principal axes; for this family the products of inertia are identically zero (measured $\max|I_{ij}|=0$ for $i\neq j$, not merely small), because it retains both mirror planes. What is compared is thus geometry at matched commanded bandwidth, not geometry plus an accidental tuning advantage.



\section{Results}
\label{sec:results}

\begin{table}[htb]
  \centering
  \caption{Maximum completing speed (\si{\metre\per\second}) by arm half-angle
  and corner sharpness. Bold marks the best in each column. Roll authority
  $\alpha_\phi$ is in \si{\radian\per\second\squared} per newton.}
  \label{tab:results}
  \begin{tabular}{r r ccc}
  \toprule
  & & \multicolumn{3}{c}{turn angle $\theta$} \\
  \cmidrule(lr){3-5}
  $t$ (\si{\degree}) & $\alpha_\phi$ & $60^\circ$ & $90^\circ$ & $120^\circ$ \\
  \midrule
    20.44 & 313.3 & 10.086 & 5.555 & 5.281 \\
    28.63 & 233.0 & 10.242 & 6.297 & 6.375 \\
    36.81 & 187.9 & \textbf{10.398} & 6.961 & 6.219 \\
    45.00 & 159.9 & 10.203 & \textbf{8.250} & 6.688 \\
    53.19 & 141.6 & 10.047 & 8.172 & \textbf{7.078} \\
    61.37 & 129.3 & 9.969 & 8.133 & \textbf{7.078} \\
    69.56 & 121.3 & 10.008 & \textbf{8.211} & \textbf{7.117} \\
  \bottomrule
\end{tabular}

\end{table}

The narrow end swaps ends. At $60^\circ$ the narrowest frame, which has the
highest roll authority, is \emph{last} (\SI{\NarrowSixty}{\percent}); at
$120^\circ$ it is \emph{first}, and the widest frame is last
(\SI{\WideOneTwenty}{\percent}). The optimum migrates monotonically with corner
sharpness: a broad plateau across $\ArgmaxSixty^\circ$ at $60^\circ$
(\NTiesSixty{} bodies within tolerance), a single peak at
$\ArgmaxNinety^\circ$ at $90^\circ$, and $\ArgmaxOneTwenty^\circ$ at
$120^\circ$. At $120^\circ$ the ranking is monotone in $t$ from best to worst,
with no inversions.

So does the size of the effect. The spread across the family grows with corner
sharpness---\SI{\SpreadSixty}{\percent} at $60^\circ$,
\SI{\SpreadNinety}{\percent} at $90^\circ$, \SI{\SpreadOneTwenty}{\percent} at
$120^\circ$---so geometry matters progressively more as the task gets harder,
which is the same conclusion read a second way.

The mechanism is the one the authorities predict. Sharper corners demand faster
bank reversals---at \SI{8.5}{\metre\per\second} on a \SI{2.4}{\metre} leg the
drone has roughly a quarter of a second per corner---which rewards roll angular
acceleration. Gentle corners flown at \SI{\BestSixty}{\metre\per\second} do not,
and there the narrow frame's weak pitch authority
($\alpha_\theta=\num{\PitchNarrow}$ against $\num{\RollNarrow}$ in roll) costs
it.

Figure~\ref{fig:flights} shows what ``completing the course'' looks like at
these speeds: the reference is tracked only loosely, and the margin at the gate
planes is small.

\paragraph{Revised 2026-08-17.} Every \emph{speed} here is a re-measurement;
the geometry and authority figures are computed from the decoder and are
unaffected. The
completion test used previously, \texttt{GateChecker}, computes lateral error
from the horizontal components of position alone, so altitude was never checked
and a drone sagging below the gates still scored passes; a second test, minimum
distance to the gate centre over the whole flight, ignores \emph{when} that
minimum occurred and so counts gates the drone passed beside rather than
through. At the $90^\circ$ operating point the three criteria score the same
flight 15/15, 15/15 and 8/15. Under the criterion defined above the best body
reaches \SI{\BestSixty}{\metre\per\second} at $60^\circ$ against
\SI{\PubBestSixty}{\metre\per\second} previously (\SI{-16.5}{\percent}),
\SI{\BestNinety}{\metre\per\second} at $90^\circ$ (\SI{-2.5}{\percent}) and
\SI{\BestOneTwenty}{\metre\per\second} at $120^\circ$ (\SI{-6.1}{\percent}).
The conclusion is unchanged in direction and cleaner in form: previously the
$120^\circ$ optimum sat mid-family, and it now sits at the narrow end with a
monotone ranking behind it. Full account in
\texttt{docs/drone\_morphology\_gate\_racing.md}~\S3.8.

\begin{figure}[htb]
  \centering
  \includegraphics[width=\textwidth]{figures/flights.pdf}
  \caption{Reference against reality, each course flown by its own best airframe
  at its own limiting speed. Grey is the B-spline the controller is given; blue
  is the path actually flown; green bars are the fifteen scored gates, all
  passed at these operating points. The drone does not track the reference---it
  lags in phase and \emph{overshoots} the weave, flying a wider amplitude than
  asked, by a mean of \SIrange{0.28}{0.33}{\metre} and a maximum of
  \SI{\XTrackNinety}{\metre}. Against the \SI{0.5}{\metre} gate radius that
  is most of the margin, which is why raising the speed a little above these
  values starts missing gates rather than degrading smoothly: the limit is set
  by cross-track overshoot at the gate planes, not by the vehicle running out of
  thrust. Note that gate passage is judged in the gate plane at the crossing,
  altitude included, so a trajectory that looks close in this projection can
  still fail.}
  \label{fig:flights}
\end{figure}

\paragraph{Heading, and the limits of the roll argument.} Yaw is commanded to
follow the trajectory tangent, which is both what the gate geometry describes
and what a vision-guided racer must do. The airframe cannot fully deliver it:
for coplanar rotors yaw comes only from rotor drag torque, never from a thrust
lever arm, so yaw authority is \SI{6.43}{\radian\per\second\squared\per\newton}
against \SIrange{121}{313}{} in roll---weaker by \numrange{19}{49}. On the
$90^\circ$ course the drone achieves \SI{34}{\degree} of yaw amplitude against
\SI{60}{\degree} commanded and crabs a mean \SI{20}{\degree}. Two consequences.
The comparison is unaffected: yaw authority is \emph{identical} across the
family---the azimuth sweep rotates rotors in-plane and leaves the drag-torque
sum untouched---so the lag is common-mode. But the mechanism above is
``mostly roll'' rather than ``roll'': with that much sideslip some lateral
acceleration is served by pitch. Roll still dominates, $\pm78^\circ$ against
$\pm28^\circ$.

\paragraph{What the speeds are, and are not.} \texttt{--speed} sets
$T_{\text{total}} = T_{\text{startup}} + \ell/v$, so $v$ is a nominal figure
rather than a speed held anywhere on the course. The median reference speed at
the gates is \SI{\GateSpeedSixty}{\metre\per\second} at $60^\circ$,
\SI{\GateSpeedNinety}{\metre\per\second} at $90^\circ$ and
\SI{\GateSpeedOneTwenty}{\metre\per\second} at $120^\circ$, against nominal
figures of \BestSixty{}, \BestNinety{} and \BestOneTwenty{}. Two independent
causes multiply to give that shortfall, and they meet at an intermediate
quantity: the mean speed over the whole path, which the first carries the
nominal figure down to and the second carries down from. The gate figure is a
median because the crossings are spread---\SIrange{\GateLoNinety}{\GateHiNinety}{\metre\per\second}
at $90^\circ$---while the whole-path figure is a mean by construction, being a
distance divided by a time.

\emph{Startup accounting.} The spline is traversed over the whole
$T_{\text{total}}$, ramp included, not over the $\ell/v$ cruise window the
formula suggests, so the mean speed along the path is $\ell_{\text{arc}} /
T_{\text{total}}$---\SI{\MeanSpeedSixty}{}, \SI{\MeanSpeedNinety}{} and
\SI{\MeanSpeedOneTwenty}{\metre\per\second} for the three courses. A second,
smaller term sits inside this one: $\ell$ is the polyline through the
waypoints, while the interpolating spline necessarily bows outside it
(\SIrange{\ArcLo}{\ArcHi}{\metre} of arc against \SI{\Polyline}{\metre}); a
curve through a fixed sequence of points is never shorter than the straight
segments joining them. That extra length is covered within the same
$T_{\text{total}}$, so it \emph{raises} the mean speed and thereby narrows the
gap to the nominal figure rather than widening it. Net of the two, the nominal
figure sits a factor of \StartupFacSixty{}, \StartupFacNinety{} and
\StartupFacOneTwenty{} above that whole-path mean.

\emph{Non-uniform parameterisation.} The spline advances uniformly in $u$
rather than by arc length, and world speed is $|dP/du|\,\dot u$
(Figure~\ref{fig:splineparam}). The clamped knot vector repeats its end knots,
so the first and last spans each carry two waypoints against one for every
interior span---\SI{\GearEnd}{\metre} of path per unit of $u$ against
\SI{\GearInterior}{\metre}, a factor of \GearRatio{}. At constant $\dot u$ the
gates are therefore geared low while the final span, which lies beyond the last
gate, sprints to about \SI{\PeakRef}{\metre\per\second} where no result is
scored. This carries the whole-path mean down to the median at the gates, a
further factor of \GearingFacSixty{}, \GearingFacNinety{} and
\GearingFacOneTwenty{}, and it is the only one of the two that grows with
corner sharpness.

\begin{figure}[htb]
  \centering
  \includegraphics[width=\textwidth]{figures/spline_parameter.pdf}
  \caption{Why the spline parameter is not arc length, on the $90^\circ$
  course. \textbf{(a)} and \textbf{(b)} show two equal intervals of $u$, each
  dotted twelve times per span and drawn to the same scale. The two end spans
  in \textbf{(a)} cover half again as much ground as the two interior spans in
  \textbf{(b)}, so their dots are spread further apart---and dot spacing is
  exactly what $\dot u$ converts into speed. The labels in \textbf{(a)} give
  the cause: the clamped knot vector puts waypoints at $u=0$, $1/3$ and
  \num{1}, three of them inside one span, against one per span in
  \textbf{(b)}. \textbf{(c)} is the same fact as a curve. The shaded end spans
  run at \SI{\GearEnd}{\metre} per unit of $u$ against
  \SI{\GearInterior}{\metre} inside, exactly a factor of \GearRatio{}; the
  ripple within each interior span is the corner itself. Every scored gate lies
  in the low-geared interior.}
  \label{fig:splineparam}
\end{figure}

The products, \SpeedRatioSixty{}, \SpeedRatioNinety{} and
\SpeedRatioOneTwenty{}, are the ratios quoted
above. Every airframe flies the identical reference, so comparisons within a
column are unaffected and bisection stays monotone. Two things follow all the
same: the numbers here are an ordinal difficulty scale rather than physical
gate speeds, and they are \textbf{not comparable across turn angles}, since
each column is divided by a different factor. All of this is measured on the
reference trajectory alone by \texttt{docs/tools/speed\_semantics.py}.

% DRAFT (open item 9, 2026-09-11) -- for the author to reword.
Re-expressed as the reference's median speed at the gates, with each airframe
at its own limiting speed, the within-column spreads quoted above shrink to
\SI{\GateSpreadSixty}{\percent}, \SI{\GateSpreadNinety}{\percent} and
\SI{\GateSpreadOneTwenty}{\percent}. They shrink because the ratio rises
slightly with $v$, through the startup term $1 + T_{\text{startup}}v/\ell$, and
so discounts the fastest airframes most. The ranking within every column is
unchanged, and the spread still grows with corner sharpness. These are the
reference's speeds at the gates, not the drone's, which lags behind it.

\paragraph{Limits.} The $60^\circ$ plateau is at resolution: \NTiesSixty{}
bodies lie within the \SI{0.0625}{\metre\per\second} bisection tolerance of the
best, so that column establishes only that both extremes are worse than the
middle. At $120^\circ$ the top \NTiesOneTwenty{} are likewise tied, though there
the whole ranking is monotone in $t$, which the tie does not disturb. The
\emph{direction} of the shift is consistent across all three columns; the exact
argument of the maximum at $60^\circ$ is not resolved. All runs
use one uniform course per turn angle and one controller tuning, optimised on the
X quad. The \emph{commanded} bandwidth is identical across bodies by
construction, but the \emph{achievable} one is not, since torque margin and
motor saturation differ with geometry; a per-airframe tuning could therefore
still move the ranking.

% DRAFT (open item 4, 2026-09-16) -- for the author to reword.
% Source: docs/data/gyroscopic_counterfactual.md and docs/data/gyroscopic_mechanism.md.
More seriously, the simulated plant omits the gyroscopic term
$\Omega\times I\Omega$ of Euler's equations. The attitude controller cancels that
term, so on this plant the cancellation injects a torque instead. MuJoCo
includes the term, and PhysX enables it by default. Restoring it shifts the
optimum to wider frames at $90^\circ$ and $120^\circ$, reversing the trend
reported here, with either of two motor-mixing schemes tested. The term's yaw
component opposes the yaw a slalom demands on narrow frames and assists it on
wide ones. On the most asymmetric frames it exceeds the yaw torque the rotors can
produce. The ranking above should therefore be read as a property of the reduced
plant until it is regenerated with the full rigid-body dynamics.

\section{Implication}

No single static geometry is best across the difficulty range. That is the first
experimental support in this work for making arm geometry adjustable in flight:
a drone that can sweep its arms could sit at the right point for the corner it is
entering. One consequence is already visible in the tuning: an airframe that
sweeps its arms independently loses the mirror planes this family relies on,
its products of inertia become non-zero, and the diagonal gain form above stops
delivering the bandwidth it is asked for---measured spreads of
\SIrange{22.4}{26.0}{\radian\per\second} against a \SI{24}{\radian\per\second}
target for single-arm perturbations. Deriving the gains from the full inertia
tensor restores it exactly; that comparison must be in place before morphologies
are optimised, or the search rewards whichever body happens to receive the most
bandwidth. Measured before the cascade was tuned, the same sweep reported that one
airframe won everywhere---an artefact of measuring where the controller, not the
airframe, was the limit.

\section*{Provenance}

Figures, the results table and every number quoted inline
(\texttt{numbers.tex}) are generated from the repository by
\texttt{Reports/make\_figures.py}: geometry from the decoder and plant interface,
speeds from the CSVs in \texttt{docs/data/}. The inline numbers are generated
because one of them drifted---the raw torque range was quoted over
\SIrange{24}{66}{\degree}, an earlier sweep's span, rather than over the
feasible family defined here. The figures are \emph{not} tracked
--- this repository gitignores \texttt{*.pdf} --- so the generator must be run
before the first build of a fresh clone. Rebuild with

\begin{quote}\ttfamily\small
uv run --no-sync python Reports/make\_figures.py \\
latexmk -pdf Reports/morphology\_gate\_racing.tex
\end{quote}

\begin{sloppypar}
Source data: \texttt{tuned\_strict\_completion.csv}, all three turn angles under
the strict gate criterion. The superseded sequential-criterion runs
(\texttt{tuned\_60deg.csv}, \texttt{tuned\_90deg.csv},
\texttt{tuned\_60\_120deg.csv}) are retained for comparison but no longer feed
this report. Commands, the
derivations behind every metric, and the defects fixed to make these
measurements possible are recorded in
\texttt{docs/drone\_morphology\_gate\_racing.md}, whose own claims are checked
against the code by \texttt{docs/tools/check\_doc\_claims.py} in the unit suite.
\end{sloppypar}


\appendix

\section{Feedforward formulation of trajectory tracking}\label{sec:feedforward}

With feedforward enabled, the Lee (geometric) controller forms its acceleration
command as

\[
  a_{\rm cmd}
  = \underbrace{K_p\,(p_{\rm ref} - p)}_{\text{position error}}
  \;+\; \underbrace{K_v\,(\dot p_{\rm ref} - \dot p)}_{\text{velocity error}}
  \;+\; \underbrace{\ddot p_{\rm ref}}_{\text{feedforward}}
\]

which is converted to force by $F = m(a_{\rm cmd} - g)$.

With it disabled, $\dot p_{\rm ref}$ defaults to zero and $\ddot p_{\rm ref}$
is dropped:

\[
  a_{\rm cmd} = K_p\,(p_{\rm ref} - p) \;-\; K_v\,\dot p .
\]
The middle term has become pure damping against motion: the controller is then
tuned to hold position at $p_{\rm ref}$ rather than to follow it, and so chases
the reference instead of tracking it.
\end{document}
